\chapter{Turbulent Topology Optimization Results}
\label{ch:benchmark_results}

\section{Introduction}
\label{sec:ch6_introduction}
Building upon the FEniCS-SA solver verified in Chapter~\ref{ch:sa_verification} and the topology optimization methodology developed in Chapter~\ref{ch:topology_optimization_methodology}, this chapter evaluates the turbulent topology optimization solver on the pipe-bend and U-bend internal-flow benchmarks from Bayat et al.~\cite{Bayat2026} and the diffuser benchmark from Yoon~\cite{Yoon2016}. Table~\ref{tab:ch6_benchmark_cases} gives the operating points and validation analyses. Each benchmark section then introduces only the geometry, objective, and interpretation needed for that case.

The optimized layouts are interpreted through projected density fields, velocity fields, and independent sharp-geometry CFD validation in Ansys Fluent. Appendix~\ref{app:C_topology_optimization_benchmarks} retains the supporting frozen-sensitivity verification, FEniCS convergence histories, numerical settings, and Ansys grid-independence studies, grouped by benchmark case.

\section{Benchmark Cases and Evaluation Metrics}
\label{sec:ch6_strategy_metrics}

The three benchmarks cover curved internal flows and a diffuser with a prescribed outlet contraction. The Bayat cases use reciprocal wall-distance penalization, topology-dependent SA and wall-distance reaction terms, an average inlet-pressure objective, and a frozen-turbulence adjoint. The diffuser uses the same frozen-SA optimization machinery, but it minimizes the dissipation objective $J_D$.

\begin{table}[!ht]
	\centering
	\footnotesize
	\caption{Benchmark cases evaluated in Chapter~\ref{ch:benchmark_results}. Detailed physical and numerical parameters are listed in Appendix~\ref{app:C_topology_optimization_benchmarks}.}
	\label{tab:ch6_benchmark_cases}
	\renewcommand{\arraystretch}{1.2}
		\begin{tabularx}{\textwidth}{l c c >{\raggedright\arraybackslash}X}
		\hline
		\textbf{Case} & \textbf{$V_f$} & \textbf{$Re$} & \textbf{Analyses reported} \\
		\hline
			Pipe bend~\cite{Bayat2026} & $0.25$ & $5{,}000$ & FEniCS optimization; Ansys-SA and Ansys SST $k-\omega$ CFD validation; laminar-design comparison \\
			U-bend~\cite{Bayat2026} & $0.27$ & $5{,}000$ & FEniCS optimization; Ansys-SA and Ansys SST $k-\omega$ CFD validation; laminar-design comparison \\
			Diffuser~\cite{Yoon2016} & $0.30$ & $3{,}000$ & Velocity-outlet baseline; pressure-outlet FEniCS optimization with $J_D$; Ansys-SA and Ansys SST $k-\omega$ CFD validation; laminar-design comparison for the pressure-outlet case \\
		\hline
	\end{tabularx}
\end{table}

The pipe-bend and U-bend cases minimize the average inlet-pressure objective $J_p$ defined in Eq.~\ref{eq:inlet_pressure_objective}, which keeps their setup comparable with the Bayat benchmark objective. The diffuser case minimizes the dissipation objective $J_D$ defined in Eq.~\ref{eq:dissipation_objective}, matching the formulation used by Yoon~\cite{Yoon2016}. All extracted geometries are then evaluated as sharp-wall Ansys Fluent models, where the Brinkman penalty is removed and the validation metrics are recomputed on the exported flow domain.

The primary validation metric is the physical viscous dissipation integrated over the full exported CFD domain, including the non-design inlet, outlet, and fixed regions that are retained in the final geometry,
\begin{equation}
	\Phi_D = \int_{\Omega}
	\frac{1}{2} \mu_{\mathrm{eff}}
	\left(\nabla\mathbf{u}+\nabla\mathbf{u}^{T}\right):
	\left(\nabla\mathbf{u}+\nabla\mathbf{u}^{T}\right)
	\,d\Omega ,
	\label{eq:ch6_validation_dissipation}
\end{equation}
where $\Omega$ denotes the complete sharp-geometry flow domain exported to Ansys and $\mathbf{S}=\frac{1}{2}\left(\nabla\mathbf{u}+\nabla\mathbf{u}^{T}\right)$ is the symmetric strain-rate tensor. The integrand in Eq.~\ref{eq:ch6_validation_dissipation} can therefore also be written as $2\mu_{\mathrm{eff}}\mathbf{S}:\mathbf{S}$. In the Ansys Fluent post-processing, the scalar strain-rate magnitude is $\dot{\gamma}=\sqrt{2\mathbf{S}:\mathbf{S}}$, so the custom field function $\mu_{\mathrm{eff}}\dot{\gamma}^{2}$ is equivalent to the integrand in Eq.~\ref{eq:ch6_validation_dissipation}. Because the reported Ansys validations all re-solve the extracted geometries as turbulent flows, the effective viscosity used in these post-processed metrics is $\mu_{\mathrm{eff}}=\mu+\mu_t$.

The second global validation metric is the static pressure drop between the physical inlet and outlet sections,
\begin{equation}
	\Delta p =
	\bar{p}_{\mathrm{in}}-\bar{p}_{\mathrm{out}},
	\qquad
	\bar{p}_{\Gamma} =
	\frac{1}{|\Gamma|}
	\int_{\Gamma} p_s\,d\Gamma ,
	\label{eq:ch6_static_pressure_drop}
\end{equation}
where $p_s$ is the static pressure and $|\Gamma|=\int_{\Gamma}1\,d\Gamma$ is the boundary length in the two-dimensional model. In Ansys Fluent, $p_s$ is taken directly from the area-weighted average static pressure on the inlet and outlet surfaces. For FEniCS results that use kinematic pressure, the corresponding static pressure is obtained as $p_s=\rho p_k$ before evaluating Eq.~\ref{eq:ch6_static_pressure_drop}. Pressure drop is retained throughout Chapter~\ref{ch:benchmark_results} because it is the most direct engineering interpretation of the same mechanical losses measured by $\Phi_D$. All two-dimensional dissipation quantities are evaluated per unit depth.

\section{Pipe Bend Benchmark}
\label{sec:alexandersen_pipebend}
The pipe-bend benchmark introduces a direct comparison with the turbulent topology optimization results of Bayat et al.~\cite{Bayat2026}. While the reference study employs a standard $k$-$\varepsilon$ RANS model with wall functions, the present framework tackles the same problem using a wall-resolved SA model under the frozen-turbulence assumption. This comparison tests whether the fundamental topological solution, a streamlined and separation-avoiding bend, remains consistent across different turbulence closures.

\subsection{Geometry and Boundary Conditions}
\label{subsec:alexandersen_pipebend_setup}
The boundary conditions and geometric constraints replicate the setup shown in Fig.~10 of Bayat et al.~\cite{Bayat2026}, specifically targeting the $V_f=0.25$ volume-fraction case. The inlet is placed on the left side of the non-design extension, the outlet is placed along the lower boundary, and the outlet pressure is fixed to zero. The optimization minimizes the average pressure on this prescribed inlet boundary required to drive the flow through the evolving bend.
The local benchmark geometry used for the present implementation is shown in Figure~\ref{fig:alexandersen_pipebend_setup}.

\begin{figure}[!ht]
	\centering
	\includegraphics[width=0.85\textwidth]{PipeBend-Geometry-BCs}
	\caption{Pipe-bend benchmark geometry and boundary conditions, adapted from Fig.~10 of Bayat et al.~\cite{Bayat2026} \protect\ChatGPTFigureRef.}
	\label{fig:alexandersen_pipebend_setup}
\end{figure}

\subsection{Optimization Results}
\label{subsec:alexandersen_pipebend_fenics}
Figure~\ref{fig:alexandersen_pipebend_fenics_results} illustrates the projected density field $\bar{\gamma}$ and the corresponding velocity magnitude. The supporting coordinate finite-difference and Taylor checks for the frozen adjoint are reported in Appendix~\ref{app:C_topology_optimization_benchmarks}. The FEniCS-SA result shows the same broad routing strategy as the $k$-$\varepsilon$ reference from Bayat et al.~\cite{Bayat2026}: the main high-speed path bypasses the sharpest corners to reduce momentum loss. A direct overlay is provided in Figure~\ref{fig:app_C_alexandersen_pipebend_paper_overlay}. The exact gray-scale density distribution near the walls still differs because the turbulence closure and near-wall treatment are not the same.

\begin{figure}[!ht]
	\centering
	\begin{subfigure}[b]{0.47\textwidth}
		\centering
		\includegraphics[width=\textwidth]{PipeBend-Turbulent-Topology-Jp}
		\caption{Projected density field $\bar{\gamma}$.}
		\label{fig:alexandersen_pipebend_topology}
	\end{subfigure}
	\hfill
	\begin{subfigure}[b]{0.49\textwidth}
		\centering
		\includegraphics[width=\textwidth]{PipeBend-Turbulent-Velocity-FEniCS-Jp}
		\caption{Velocity magnitude and vector field.}
		\label{fig:alexandersen_pipebend_velocity}
	\end{subfigure}
	\caption{FEniCS-SA optimization result for the pipe-bend benchmark at $V_f=0.25$. Final objective value $J_p= 0.2130$ Pa.}
	\label{fig:alexandersen_pipebend_fenics_results}
\end{figure}

For the same $V_f=0.25$, $Re=5{,}000$ operating point, Bayat et al.~\cite{Bayat2026} report $J_p=0.4023$ Pa for their conventional $k$-$\varepsilon$ topology optimization. This value is used only as a reference scale, not as a direct validation target, because the turbulence closure and near-wall treatment differ from the present wall-resolved SA formulation.

The reported FEniCS value is the endpoint of a short $70$-iteration run without continuation. It is used here as a provisional design state for comparing topology and validation metrics, not as a converged optimum. The convergence history in Figure~\ref{fig:alexandersen_pipebend_histories} shows that the normalized pressure objective is still monotonically decreasing, so the run is promising but should be extended with more MMA iterations and continuation.


\subsection{CFD Validation}
\label{subsec:alexandersen_pipebend_cfd}
The extracted $70$-iteration pipe-bend topology is evaluated in Table~\ref{tab:alexandersen_pipebend_validation_metrics}. Because the SST $k-\omega$ model is notably more sensitive to adverse pressure gradients than SA, evaluating the extracted geometry under both closures serves as a stress test for the provisional shape.

\begin{table}[!ht]
	\centering
	\footnotesize
	\caption{CFD validation metrics for the extracted pipe-bend topology.}
	\label{tab:alexandersen_pipebend_validation_metrics}
	\renewcommand{\arraystretch}{1.2}
	\begin{tabularx}{\textwidth}{l c c c X}
		\hline
		\textbf{Solver} & \textbf{Model} & \textbf{$\Phi_D$ [W/m]} & \textbf{$\Delta p$ [Pa]} & \textbf{Comment} \\
		\hline
		FEniCS diffuse design & SA & 0.0138 & 0.2130 & 70-iteration non-continuation state \\
		Ansys sharp geometry & SA & 0.0153 & 0.1457 & Cross-code comparison \\
		Ansys sharp geometry & SST $k-\omega$ & 0.0166 & 0.1522 & Turbulence-model sensitivity \\
		\hline
	\end{tabularx}
\end{table}

Using the Ansys-SA result as the reference, the SST $k-\omega$ evaluation increases the pressure drop by approximately $4.5\%$ and the viscous dissipation by approximately $8.5\%$. The SST $k-\omega$ model therefore predicts slightly higher losses, consistent with a closure that is more sensitive to adverse pressure gradients, but the changes remain modest. This suggests that the extracted pipe-bend geometry is not narrowly tuned to the SA model. It does not, however, establish that the underlying FEniCS optimization has converged.

\subsection{Comparison with Laminar Design}
\label{subsec:alexandersen_pipebend_laminar}
The laminar reference uses the same pipe-bend mesh, design-domain restrictions, inlet and outlet treatment, objective type, and volume fraction, but replaces the SA turbulent state equations by the laminar flow formulation and is optimized at $Re=1$. This laminar optimization favors a sharper, shorter turn that is appropriate for the creeping-flow reference problem but less suitable once the extracted geometry is operated at the target turbulent Reynolds number. The comparison is therefore interpreted under identical turbulent operating conditions, rather than by comparing the laminar and turbulent optimization objectives directly.

Figure~\ref{fig:alexandersen_pipebend_laminar_topologies} compares the laminar and turbulent optimized topologies before sharp-geometry extraction. The laminar-design mesh-independence study supporting the SST $k-\omega$ result is reported in Table~\ref{tab:app_C_alexandersen_pipebend_lam_mesh} of Appendix~\ref{app:C_topology_optimization_benchmarks}.

\begin{figure}[!ht]
	\centering
	\begin{subfigure}[b]{0.47\textwidth}
		\centering
		\includegraphics[width=\textwidth]{PipeBend-Laminar-Topology-Jp}
		\caption{Projected density field $\bar{\gamma}$: laminar-optimized design ($Re=1$).}
		\label{fig:alexandersen_pipebend_laminar_topo}
	\end{subfigure}
	\hfill
	\begin{subfigure}[b]{0.47\textwidth}
		\centering
		\includegraphics[width=\textwidth]{PipeBend-Turbulent-Topology-Jp}
		\caption{Projected density field $\bar{\gamma}$: turbulent-optimized design ($Re=5{,}000$).}
		\label{fig:alexandersen_pipebend_turbulent_topo}
	\end{subfigure}
	\caption{Laminar and turbulent topology comparison for the pipe-bend benchmark.}
	\label{fig:alexandersen_pipebend_laminar_topologies}
\end{figure}


Figure~\ref{fig:alexandersen_pipebend_laminar_velocity_comparison} shows the corresponding SST $k-\omega$ velocity contours through the extracted laminar and turbulent pipe-bend designs. These contours provide the flow-field evidence behind the scalar metrics: sharper turns and corners in the laminar design promote stronger local separation and momentum loss when the geometry is operated at the turbulent Reynolds number.

\begin{figure}[!ht]
	\centering
	\begin{subfigure}[b]{0.47\textwidth}
		\centering
		\includegraphics[width=\textwidth]{PipeBend-Laminar-VelocityMagnitude-SSTkomega}
		\caption{Velocity magnitude field $u$: laminar-optimized design.}
		\label{fig:alexandersen_pipebend_laminar_velocity_sst}
	\end{subfigure}
	\hfill
	\begin{subfigure}[b]{0.47\textwidth}
		\centering
		\includegraphics[width=\textwidth]{PipeBend-Turbulent-VelocityMagnitude-SSTkomega}
		\caption{Velocity magnitude field $u$: turbulent-optimized design.}
		\label{fig:alexandersen_pipebend_turbulent_velocity_sst}
	\end{subfigure}
	\caption{Velocity-contour comparison for the extracted pipe-bend designs under identical turbulent operating conditions (SST $k-\omega$ model).}
	\label{fig:alexandersen_pipebend_laminar_velocity_comparison}
\end{figure}

The final SST $k-\omega$ comparison in Table~\ref{tab:alexandersen_pipebend_laminar_perf} evaluates the extracted laminar and turbulent pipe-bend designs under identical turbulent operating conditions. The turbulent design reduces the viscous dissipation from $0.0313$ to $0.0166~\mathrm{W/m}$ and the pressure drop from $0.5444$ to $0.1522~\mathrm{Pa}$. These reductions, $47.0\%$ and $72.0\%$ respectively, show that the turbulent topology optimization framework outperforms the laminar surrogate for the target turbulent operating point.

\begin{table}[!ht]
	\centering
	\footnotesize
	\caption{Performance comparison of the extracted pipe-bend topologies under identical turbulent operating conditions.}
	\label{tab:alexandersen_pipebend_laminar_perf}
	\renewcommand{\arraystretch}{1.2}
	\begin{tabularx}{\textwidth}{X c c c X}
		\hline
		\textbf{Design origin} & \textbf{Model} & \textbf{$\Phi_D$ [W/m]} & \textbf{$\Delta p$ [Pa]} & \textbf{Relative performance} \\
		\hline
		Laminar TO, $Re=1$ & SST $k-\omega$ & $0.0313$ & $0.5444$ & Baseline \\
		Turbulent TO, SA & SST $k-\omega$ & $0.0166$ & $0.1522$ & $47.0\%$ lower $\Phi_D$; $72.0\%$ lower $\Delta p$ \\
		\hline
	\end{tabularx}
\end{table}

\section{U-Bend Benchmark}
\label{sec:alexandersen_ubend}
The U-bend is the companion benchmark to the pipe bend. It uses the same frozen-SA topology optimization framework, but the geometry forces a $180^\circ$ return path between two ports on the same side of the domain. This makes the case useful for checking whether the optimizer can form a smooth turbulent return duct without relying on the shorter $90^\circ$ routing available in the pipe-bend case.

\subsection{Geometry and Boundary Conditions}
\label{subsec:alexandersen_ubend_setup}
The geometry and boundary conditions follow Fig.~11 and Table~4 of Bayat et al.~\cite{Bayat2026}. The square design region has side length $L$, with a non-design lead duct of length $0.2L$ on the left. The inlet and outlet ports both have height $0.2L$; the top port acts as the inlet and the lower port acts as the pressure outlet. A fixed internal separator bar of thickness $0.10L$ extends from the lead duct into the design domain, forcing the flow to turn around its rounded tip. The target volume fraction is $V_f=0.27$, and the Reynolds number is $Re=5{,}000$ based on the port height and maximum inlet velocity.
Figure~\ref{fig:alexandersen_ubend_setup} shows the corresponding local geometry and boundary-condition layout used for this benchmark.

\begin{figure}[!ht]
	\centering
	\includegraphics[width=0.85\textwidth]{UBend-TO-Geometry-BCs}
	\caption{Geometry and boundary conditions for the U-bend benchmark, adapted from Fig.~11 and Table~4 of Bayat et al.~\cite{Bayat2026} \protect\ChatGPTFigureRef.}
	\label{fig:alexandersen_ubend_setup}
\end{figure}

\subsection{Optimization Results}
\label{subsec:alexandersen_ubend_fenics}
Figure~\ref{fig:alexandersen_ubend_fenics_results} reports the final projected density and velocity-magnitude field for the U-bend. The main feature to assess is whether the design forms a rounded return passage around the separator tip, rather than a short path with tight local curvature. Because the objective is average inlet pressure, a successful layout reduces the pressure required to maintain the imposed inlet flow while respecting the prescribed volume fraction.

\begin{figure}[!ht]
	\centering
	\begin{subfigure}[b]{0.45\textwidth}
		\centering
		\includegraphics[width=\textwidth]{UBend-Turbulent-Topology-Jp}
		\caption{Projected density field $\bar{\gamma}$.}
		\label{fig:alexandersen_ubend_topology}
	\end{subfigure}
	\hfill
	\begin{subfigure}[b]{0.53\textwidth}
		\centering
		\includegraphics[width=\textwidth]{UBend-Turbulent-Velocity-FEniCS-Jp}
		\caption{Velocity magnitude and vector field.}
		\label{fig:alexandersen_ubend_velocity}
	\end{subfigure}
	\caption{FEniCS-SA optimization result for the U-bend benchmark at $V_f=0.27$. Final objective value $J_p=0.8250~\mathrm{Pa}$.}
	\label{fig:alexandersen_ubend_fenics_results}
\end{figure}

For the matching $V_f=0.27$, $Re=5{,}000$ U-bend case, the conventional $k$-$\varepsilon$ result reported by Bayat et al.~\cite{Bayat2026} is $J_p=0.8327$. As for the pipe bend, this value provides a reference scale only; the comparison remains qualitative because the present solver uses a different turbulence model and near-wall treatment.

The turbulent U-bend topology should also be read in the same way as the provisional pipe-bend topology: it is an extracted optimization endpoint from the available run, not a fully converged final optimum. The available U-bend convergence history is reported in Figure~\ref{fig:alexandersen_ubend_histories} of Appendix~\ref{app:C_topology_optimization_benchmarks}, using the same combined objective and volume-residual format as the pipe-bend diagnostic.

\subsection{CFD Validation}
\label{subsec:alexandersen_ubend_cfd}
The extracted U-bend topology is validated with the same sharp-geometry workflow as the pipe bend. The Ansys-SA solve checks the design under the same closure family used in the optimization loop, while the Ansys SST $k-\omega$ solve tests whether the extracted turbulent design remains robust under a two-equation model that is more sensitive to separation and adverse pressure gradients. Table~\ref{tab:alexandersen_ubend_validation_metrics} summarizes the FEniCS-SA optimization diagnostic and the two sharp-geometry Ansys validation solves for the turbulent U-bend design; the supporting validation mesh hierarchy is documented in Table~\ref{tab:app_B_alexandersen_ubend_turb_mesh} of Appendix~\ref{app:C_topology_optimization_benchmarks}.

\begin{table}[!ht]
	\centering
	\footnotesize
	\caption{CFD validation metrics for the extracted turbulent U-bend topology.}
	\label{tab:alexandersen_ubend_validation_metrics}
	\renewcommand{\arraystretch}{1.2}
	\begin{tabularx}{\textwidth}{l c c c X}
		\hline
		\textbf{Solver} & \textbf{Model} & \textbf{$\Phi_D$ [W/m]} & \textbf{$\Delta p$ [Pa]} & \textbf{Comment} \\
		\hline
		FEniCS diffuse design & SA & $0.0936$ & $0.8250$ & Optimization-state diagnostic \\
		Ansys sharp geometry & SA & $0.0785$ & $0.8071$ & Cross-code comparison with the optimization closure family \\
		Ansys sharp geometry & SST $k-\omega$ & $0.1016$ & $1.1973$ & Independent turbulence-closure check \\
		\hline
	\end{tabularx}
\end{table}

Using the Ansys-SA result as the reference, the SST $k-\omega$ validation predicts larger losses for the same extracted turbulent geometry: $\Phi_D$ increases by approximately $29.4\%$ and $\Delta p$ by approximately $48.3\%$. This trend is physically plausible because the fixed separator forces a compact return flow, where SST $k-\omega$ is more responsive to adverse pressure gradients and local recirculation downstream of the separator. The U-bend validation is therefore interpreted as a closure-sensitivity result for the forced $180^\circ$ return flow rather than as closure-independent agreement.

\subsection{Comparison with Laminar Design}
\label{subsec:alexandersen_ubend_laminar}
The laminar reference uses the same U-bend mesh, inlet and outlet ports, separator geometry, objective type, and volume fraction, but it is optimized at $Re=1$ without the SA turbulence model. The extracted laminar and turbulent layouts are then evaluated under the same target turbulent operating conditions. This comparison tests whether the SA-optimized turbulent return path remains advantageous when the sharp geometry is re-solved with an independent turbulence closure.

Figure~\ref{fig:alexandersen_ubend_laminar_topologies} compares the laminar and turbulent optimized U-bend topologies. The laminar-design mesh-independence study supporting the SST $k-\omega$ result is reported in Table~\ref{tab:app_C_alexandersen_ubend_lam_mesh} of Appendix~\ref{app:C_topology_optimization_benchmarks}.

\begin{figure}[!ht]
	\centering
	\begin{subfigure}[b]{0.47\textwidth}
		\centering
		\includegraphics[width=\textwidth]{UBend-Laminar-Topology-Jp}
		\caption{Projected density field $\bar{\gamma}$: laminar-optimized design ($Re=1$).}
		\label{fig:alexandersen_ubend_laminar_topo}
	\end{subfigure}
	\hfill
	\begin{subfigure}[b]{0.50\textwidth}
		\centering
		\includegraphics[width=\textwidth]{UBend-Turbulent-Topology-Jp}
		\caption{Projected density field $\bar{\gamma}$: turbulent-optimized design ($Re=5{,}000$).}
		\label{fig:alexandersen_ubend_turbulent_topo}
	\end{subfigure}
	\caption{Laminar and turbulent topology comparison for the U-bend benchmark.}
	\label{fig:alexandersen_ubend_laminar_topologies}
\end{figure}

Figure~\ref{fig:alexandersen_ubend_laminar_velocity_comparison} shows SST $k-\omega$ velocity contours through the extracted U-bend designs, plotted with the same color scale. The fields show broadly similar return-flow organization, with a high-speed outer path around the bend and a low-momentum band along the inner side downstream of the separator. The close-up vector diagnostic in Figure~\ref{fig:app_C_alexandersen_ubend_recirc_sst} of Appendix~\ref{app:C_topology_optimization_benchmarks} shows local recirculation in this low-momentum region, so the velocity fields are interpreted as qualitative evidence of curvature-induced separation and recovery rather than as the primary performance metric.

\begin{figure}[!ht]
	\centering
	\begin{subfigure}[b]{0.47\textwidth}
		\centering
		\includegraphics[width=\textwidth]{UBend-Laminar-VelocityMagnitude-SSTkomega}
		\caption{Velocity magnitude field $u$: laminar-optimized design.}
		\label{fig:alexandersen_ubend_laminar_velocity_sst}
	\end{subfigure}
	\hfill
	\begin{subfigure}[b]{0.47\textwidth}
		\centering
		\includegraphics[width=\textwidth]{UBend-Turbulent-VelocityMagnitude-SSTkomega}
		\caption{Velocity magnitude field $u$: turbulent-optimized design.}
		\label{fig:alexandersen_ubend_turbulent_velocity_sst}
	\end{subfigure}
	\caption{SST $k-\omega$ velocity-contour comparison for the extracted U-bend designs under identical turbulent operating conditions.}
	\label{fig:alexandersen_ubend_laminar_velocity_comparison}
\end{figure}

Table~\ref{tab:alexandersen_ubend_laminar_perf} gives the corresponding laminar-design comparison under identical turbulent operating conditions.

\begin{table}[!ht]
	\centering
	\footnotesize
	\caption{Performance comparison of the extracted U-bend topologies under identical turbulent operating conditions.}
	\label{tab:alexandersen_ubend_laminar_perf}
	\renewcommand{\arraystretch}{1.2}
	\begin{tabularx}{\textwidth}{X c c c X}
		\hline
		\textbf{Design origin} & \textbf{Validation model} & \textbf{$\Phi_D$ [W/m]} & \textbf{$\Delta p$ [Pa]} & \textbf{Relative performance} \\
		\hline
		Laminar TO, $Re=1$ & SA & $0.1137$ & $0.9646$ & Baseline under SA \\
		Turbulent TO, SA & SA & $0.0785$ & $0.8071$ & $31.0\%$ lower $\Phi_D$; $16.3\%$ lower $\Delta p$ \\
		Laminar TO, $Re=1$ & SST $k-\omega$ & $0.1069$ & $1.1496$ & Baseline under SST $k-\omega$ \\
		Turbulent TO, SA & SST $k-\omega$ & $0.1016$ & $1.1973$ & $5.0\%$ lower $\Phi_D$; $4.1\%$ higher $\Delta p$ \\
		\hline
	\end{tabularx}
\end{table}

Under Ansys-SA, the turbulent U-bend design gives the expected improvement over the laminar reference. Under SST $k-\omega$, however, the two designs are much closer: the turbulent design has a slightly lower viscous dissipation but a slightly higher pressure drop. The U-bend therefore supports the same cautious message as the velocity fields: the SA-optimized return path is beneficial under the design closure, but the $180^\circ$ turn remains sensitive to the turbulence model used for sharp-geometry validation.

\section{Diffuser Benchmark}
\label{sec:yoon_diffuser}
The diffuser benchmark follows the third numerical example of Yoon~\cite{Yoon2016}. It adds a downstream-boundary-condition sensitivity to the closure-sensitivity checks already used for the bend cases. A velocity-outlet baseline follows the published setup, while the primary result replaces the prescribed outlet velocity profile with a pressure outlet so that the outlet velocity distribution is solved rather than imposed.

\subsection{Velocity-Outlet Boundary-Condition Sensitivity}
\label{subsec:yoon_diffuser_velocity_outlet}
The velocity-outlet baseline prescribes parabolic velocity profiles at both the full-height inlet and the center-third outlet. It produces a more strongly diverging channel near the outlet, which is interpreted as partly boundary-condition-driven. Its geometry, FEniCS fields, and Ansys diagnostic contour are reported in Appendix~\ref{app:C_topology_optimization_benchmarks}; the pressure-outlet variant is used as the primary diffuser result below.

\subsection{Pressure-Outlet Setup and Optimization Results}
\label{subsec:yoon_diffuser_pressure_setup}
\label{subsec:yoon_diffuser_pressure_fenics}
The pressure-outlet variant changes only the downstream boundary treatment relative to the velocity-outlet baseline documented in Appendix~\ref{app:C_topology_optimization_benchmarks}. Its local geometry and boundary-condition layout are shown in Figure~\ref{fig:yoon_diffuser_pressure_setup}. The inlet remains the shifted parabolic velocity profile with $U_{\max,\mathrm{in}}=3~\mathrm{m/s}$ and the same $Re=3{,}000$ definition, while the center outlet is assigned a fixed static pressure $p=0~\mathrm{Pa}$. No outlet velocity profile is prescribed. This makes the outlet mass flux and velocity distribution part of the solution and removes the direct Dirichlet forcing that affected the velocity-outlet topology.

\begin{figure}[!ht]
	\centering
	\includegraphics[width=0.95\textwidth]{Diffuser-Geometry-BCs-PreOut}
	\caption{Pressure-outlet diffuser geometry and boundary conditions \protect\ChatGPTFigureRef.}
	\label{fig:yoon_diffuser_pressure_setup}
\end{figure}

Figure~\ref{fig:yoon_diffuser_pressure_fenics_results} reports the pressure-outlet FEniCS-SA topology and velocity field. Compared with the velocity-outlet result in Appendix~\ref{app:C_topology_optimization_benchmarks}, the outlet region is less strongly widened. The optimized pressure-outlet objective is also lower, decreasing from $J_D=3.1164\times10^5~\mathrm{W/m}$ to $J_D=2.7664\times10^5~\mathrm{W/m}$, or by about $11.2\%$. This is not a direct performance ranking between equivalent designs, because the boundary-value problem has changed. The conclusion is narrower: removing the outlet Dirichlet velocity profile reduces the amount of geometry change required near the outlet and gives a less forced diffuser design.

\begin{figure}[!ht]
	\centering
	\begin{subfigure}[b]{0.47\textwidth}
		\centering
		\includegraphics[width=\textwidth]{Diffuser-Turbulent-Topology-PreOut}
		\caption{Projected density field $\bar{\gamma}$.}
		\label{fig:yoon_diffuser_pressure_topology}
	\end{subfigure}
	\hfill
	\begin{subfigure}[b]{0.50\textwidth}
		\centering
		\includegraphics[width=\textwidth]{Diffuser-Turbulent-Velocity-FEniCS-PreOut}
		\caption{Velocity magnitude and vector field.}
		\label{fig:yoon_diffuser_pressure_velocity}
	\end{subfigure}
	\caption{FEniCS-SA optimization result for the pressure-outlet diffuser at $V_f=0.30$. Final objective value $J_D=2.7664\times10^5~\mathrm{W/m}$.}
	\label{fig:yoon_diffuser_pressure_fenics_results}
\end{figure}

The corresponding pressure-outlet convergence history through this common comparison point is reported in Figure~\ref{fig:yoon_diffuser_pressure_convergence} of Appendix~\ref{app:C_topology_optimization_benchmarks}.

\subsection{Pressure-Outlet CFD Validation}
\label{subsec:yoon_diffuser_pressure_cfd}
The extracted pressure-outlet topology is validated with the same sharp-geometry workflow as the pipe bend and U-bend. The Ansys-SA solve checks the transfer from the diffuse FEniCS design to a sharp-wall geometry under the same closure family, while the Ansys SST $k-\omega$ solve tests whether the ranking is robust under an independent two-equation model. Because the optimization objective is $J_D$, the validation emphasizes both physical viscous dissipation and static pressure drop.

\begin{table}[!ht]
	\centering
	\footnotesize
	\caption{CFD validation metrics for the extracted pressure-outlet diffuser topology.}
	\label{tab:yoon_diffuser_validation_metrics}
	\renewcommand{\arraystretch}{1.2}
		\begin{tabularx}{\textwidth}{l c c >{\raggedright\arraybackslash}X}
		\hline
		\textbf{Evaluation} & \textbf{$\Phi_D$ [W/m]} & \textbf{$\Delta p_s$ [Pa]} & \textbf{Interpretation} \\
		\hline
		FEniCS-SA diffuse topology & $2.6307\times10^5$ & $2.6022\times10^5$ & Diffuse-domain diagnostic; total objective $J_D=2.7664\times10^5~\mathrm{W/m}$ \\
		Ansys-SA sharp geometry & $1.1885\times10^5$ & $1.1541\times10^5$ & Sharp-wall validation with the optimization closure family \\
		Ansys SST $k-\omega$ sharp geometry & $1.2834\times10^4$ & $5.0006\times10^4$ & Independent closure check; much lower modeled dissipation \\
		\hline
	\end{tabularx}
\end{table}

The FEniCS-SA values in Table~\ref{tab:yoon_diffuser_validation_metrics} are approximately $2.21$ times larger in $\Phi_D$ and $2.25$ times larger in $\Delta p_s$ than the Ansys-SA sharp-geometry values, so this case cannot be claimed as quantitative cross-code agreement. The result is interpreted as a transfer limitation between the diffuse Brinkman-domain optimization diagnostic and the sharp-wall extracted geometry.

The Ansys validation rows are also strongly closure-sensitive: for the turbulent pressure-outlet design, the SA pressure drop is $2.31$ times the SST $k-\omega$ value, and the SA dissipation is $9.26$ times the SST $k-\omega$ value. This is the opposite trend from the U-bend: here SST $k-\omega$ predicts lower modeled losses, while SA appears more loss-producing in the adverse-pressure-gradient diffuser flow. Because the split is large, it is treated as model-form uncertainty rather than evidence that either closure gives a closure-independent physical truth. Additional velocity-magnitude, turbulent-viscosity-ratio, and local dissipation-density fields were inspected but are not retained as main figures because they do not change the interpretation. The appendix wall-shear diagnostic and SST $k-\omega$ vector close-up in Figures~\ref{fig:app_C_yoon_diffuser_wall_shear_sign_map} and~\ref{fig:app_C_yoon_diffuser_separation_sst} identify local reverse-flow regions in the pressure-outlet geometry; the integrated pressure-drop and dissipation metrics remain the primary evidence for this diffuser.

\subsection{Pressure-Outlet Laminar-Design Comparison}
\label{subsec:yoon_diffuser_pressure_laminar}
The laminar reference for the final diffuser comparison uses the same pressure-outlet boundary treatment, objective type, and volume fraction, but it is optimized at $Re=1$ without the SA turbulence model. The extracted laminar and turbulent pressure-outlet layouts are then evaluated under the same target turbulent operating conditions. This comparison isolates whether the SA-optimized pressure-outlet design remains beneficial when the final sharp geometry is re-solved with the same closure and with SST $k-\omega$.

Table~\ref{tab:yoon_diffuser_pressure_laminar_perf} compares the extracted laminar and turbulent pressure-outlet diffuser topologies under identical turbulent operating conditions. The topology and velocity fields were inspected, but they do not by themselves reveal the performance ranking clearly enough to justify a separate figure. Because the diffuser is optimized with $J_D$, lower $\Phi_D$ is the primary performance indicator, while the static pressure drop remains the most direct hydraulic interpretation of the extracted sharp geometry. Both SA and SST $k-\omega$ validation solves are included to quantify whether the final ranking changes with the validation model.

\begin{table}[!ht]
	\centering
	\footnotesize
	\caption{Performance comparison of the extracted pressure-outlet diffuser topologies under identical turbulent operating conditions.}
	\label{tab:yoon_diffuser_pressure_laminar_perf}
	\renewcommand{\arraystretch}{1.2}
		\begin{tabularx}{\textwidth}{l c c >{\raggedright\arraybackslash}X}
		\hline
		\textbf{Evaluation} & \textbf{$\Phi_D$ [W/m]} & \textbf{$\Delta p_s$ [Pa]} & \textbf{Relative performance} \\
		\hline
		Laminar TO, Ansys-SA & $1.3495\times10^5$ & $1.3979\times10^5$ & Baseline under SA \\
		Turbulent TO, Ansys-SA & $1.1885\times10^5$ & $1.1541\times10^5$ & $11.9\%$ lower $\Phi_D$; $17.4\%$ lower $\Delta p_s$ \\
		Laminar TO, Ansys SST $k-\omega$ & $1.1289\times10^4$ & $3.8332\times10^4$ & Baseline under SST $k-\omega$ \\
		Turbulent TO, Ansys SST $k-\omega$ & $1.2834\times10^4$ & $5.0006\times10^4$ & $13.7\%$ higher $\Phi_D$; $30.5\%$ higher $\Delta p_s$ \\
		\hline
	\end{tabularx}
\end{table}

The laminar-design comparison reinforces this closure-sensitive interpretation. Under Ansys-SA, the turbulent topology reduces the pressure drop by $17.4\%$ and the viscous dissipation by $11.9\%$ relative to the pressure-outlet laminar topology. Under SST $k-\omega$, the ranking reverses: the turbulent topology gives a $30.5\%$ higher pressure drop and a $13.7\%$ higher dissipation. The appendix wall-shear and vector diagnostics identify where local reverse flow appears in the pressure-outlet geometry, but the case-specific conclusion comes from the scalar ranking reversal. The diffuser is therefore retained as a closure-sensitivity case for the present SA-based framework.

\section{Conclusion}
\label{sec:ch6_conclusion}
This chapter answers the third research question by evaluating whether the turbulent topology optimization framework produces designs that remain credible after independent CFD validation. The pipe bend provides the strongest validation evidence among the reported cases: the Appendix~\ref{app:C_topology_optimization_benchmarks} frozen-sensitivity checks confirm the derivative supplied to MMA for $J_p$, and the extracted $70$-iteration topology remains robust when validated in Ansys Fluent with both SA and SST $k-\omega$. The FEniCS optimization endpoint itself is provisional rather than converged, but the extracted geometry is not narrowly tuned to the SA model.

The U-bend extends the assessment to a forced $180^\circ$ return path and gives a more cautious result. Like the pipe-bend topology, the turbulent U-bend topology is an optimization endpoint rather than a fully converged final optimum, so the result is interpreted as validation of the extracted geometry and not as a convergence claim. For the turbulent topology alone, replacing Ansys-SA by SST $k-\omega$ increases $\Phi_D$ by $29.4\%$ and $\Delta p$ by $48.3\%$. Compared with the laminar reference, the turbulent design improves both metrics under SA, but under SST $k-\omega$ it gives only a $5.0\%$ lower dissipation and a $4.1\%$ higher pressure drop. The appendix close-up vector diagnostic explains this case-specific sensitivity through local recirculation in the low-momentum inner return region.

The diffuser adds a separate caution: outlet treatment and validation closure both affect the final ranking. The velocity-outlet baseline shows that a prescribed outlet profile can influence the optimized outlet geometry, so the pressure-outlet variant is used as the primary diffuser result. For this pressure-outlet diffuser, the turbulent SA design improves on the laminar design under Ansys-SA, reducing $\Delta p_s$ by $17.4\%$ and $\Phi_D$ by $11.9\%$. The same two geometries evaluated with SST $k-\omega$ give the opposite ranking, with the turbulent design increasing $\Delta p_s$ by $30.5\%$ and $\Phi_D$ by $13.7\%$. Overall, the pipe bend provides the cleanest transfer case, while the U-bend and diffuser show why extracted turbulent designs must still be checked under an independent closure before making closure-independent claims.
